\documentclass[12pt]{article}
\usepackage[utf8]{inputenc}
\usepackage[T1]{fontenc}
\usepackage{hyperref}
\title{Loup B Language Resources for K-12 Curriculum Development}
\author{Compiled by Claude Sonnet 5.5 (claude-sonnet-5-5) for Andr\'as Kornai}
\date{}
\begin{document}
\maketitle

\section{Introduction}
Loup B (ISO 639-3: xlb; Glottocode: loup1245) is an extinct Algic (Eastern Algonquian) variety from colonial New England. It is known only from mid-eighteenth-century manuscripts by French missionaries. Loup B is the variety described by Magon de Terlaye, while Loup A is the one described by Mathevet. The documentation is very thin, reportedly about 14 pages. Its classification is uncertain. One account describes it as close to Mahican and Western Abenaki, and possibly a composite of dialects. The French term ``Loup'' was used inconsistently.

\section{Primary Educational Resources}
I found no primers, textbooks or school programmes for Loup B. There are no living speakers, so any heritage curriculum could only draw on the manuscript record and on related, better-documented languages.

\section{Linguistic and Reference Materials}
\begin{itemize}
\item Loup languages (Wikipedia). \url{https://en.wikipedia.org/wiki/Loup_languages}
\item Eastern Algonquian languages (Wikipedia). \url{https://en.wikipedia.org/wiki/Eastern_Algonquian_languages}
\item Glottolog languoid page from the search (listed under code xpq). It may be a related Loup entry, so check it before relying on it. \url{https://glottolog.org/resource/languoid/iso/xpq}
\end{itemize}

\section{Religious and Community Resources}
I found none specific to Loup B.

\section{Digital and Audio Resources}
I found no audio or digital corpus for Loup B.

\section{Curriculum Development Context}
Loup B is an extinct language with a very small written record. A realistic revitalization or heritage programme would focus on history, place names and the cultural context of the Pocumtuck and nearby peoples. It could use better-documented relatives such as Western Abenaki or Mahican, with the community descendants as partners. I did not locate any current community programme.

\end{document}